\chapter{平移算子和生成元}

\section{空间平移}

\pondering{量子力学里，我们平移了什么？为什么要平移？}\label{pondering:meaning}

我们知道，量子态可以理解为系统物理状态的表征，引入之的目的是为了理解测量问题.

现在考虑一个一维空间中的粒子，其抽象态为 \(|\psi\rangle\). 现在，我们关心它在哪儿，即「位置」物理量的测量值. 那么，把 \(|\psi\rangle\) 往 \(X\) 的本征态 \(|x\rangle\) 上投影，得到位置表象下的波函数
\[
\psi(x)=\langle x|\psi\rangle
\]

如果我们用一个覆盖任意区间 \([u,v]\) 的固定探测器去「捕捉」这个粒子，那测得它的概率即为粒子落在其中的概率\side{固定的探测区间}{\begin{tikzpicture}[x=0.78cm,y=0.67cm]
 \fill[accent!9] (1.3,0) rectangle (3.4,1.55);
 \draw[->,gray] (0,0)--(5,0) node[below]{$x$};
 \draw[muted,dashed,domain=0:4.7,samples=60] plot (\x,{exp(-2*(\x-1.65)^2)});
 \draw[accent,domain=0:4.7,samples=60] plot (\x,{exp(-2*(\x-2.7)^2)});
 \draw[->] (1.65,1.19)--node[above]{$a$}(2.7,1.19);
 \node[below] at (1.3,0){$u$};\node[below] at (3.4,0){$v$};
 \end{tikzpicture}%
\par 虚线为平移前的概率密度，实线为平移后的概率密度. }
\[
P_0=\int_u^v |\psi(x)|^2\,dx
\]

设整个粒子系统向右平移 \(a\)，探测器不动，记新波函数为 \(\psi_a(x)\)，现在落在 \([u,v]\) 的波包，原先落在 \([u-a,v-a]\)，有
\[
P_a=\int_u^v |\psi_a(x)|^2\,dx=\int_{u-a}^{v-a}|\psi(x)|^2\,dx
\]
变量代换：
\[
P_a=\int_u^v |\psi(x-a)|^2\,dx
\]
暂时忽略相对相位，有
\[
\psi_a(x)=\psi(x-a).
\]
这告诉我们，在空间中平移一个系统，等效于在数学上平移它的位置表象波函数.

为什么我们关心平移？因为我们关心，同样的系统，在不同的空间位置，是否有相同的演化.

\pondering{为什么要引入平移算子？}\label{pondering:operator}

显然，由 \(\int_u^v|\psi_a(x)|^2\,dx\) 和 \(\int_u^v|\psi(x-a)|^2\,dx\) 推出 \(\psi_a(x)=\psi(x-a)\)，是不严谨的，因为，一般地，\(\psi_a(x)=e^{i\theta(x,a)}\psi(x-a)\) 均符合条件. 所以接下来，我们要更严格的定义「平移」动作.

让我们先达成共识：\(x\) 是在尺子上看到的数，\(|x\rangle\) 是这个读数对应的理想定位态. \(X\) 表示「测量位置」的可观测量. \(X|x\rangle=x|x\rangle\) 即表达：若粒子精确定位在 \(x\) 处，则一定测得位置 \(x\).

现在设想我们在实验室制备了精确位于 \(x\) 位置的粒子（事实上这是不可能的），有 \(|\psi\rangle=|x\rangle\). 现在我们用某种「神秘方法」把它右移 \(a\) 到 \(x+a\) 位置，即 \(|\psi_a\rangle=|x+a\rangle\). 对这种「神秘方法」，我们定义线性平移算子 \(T(a)\) 满足
\[
T(a)|x\rangle=|x+a\rangle
\]
\side{相位约定}{注意，这事实上是一个自然且方便的定义，\(\widetilde{T}(a)=e^{ika}T(a)\) 也完全是物理的. }

同时我们知道波函数是表象对应物理量本征态的叠加系数：
\[
|\psi\rangle=\int\psi(x)|x\rangle\,dx
\]
将 \(T(a)\) 作用在两侧：
\[
T(a)|\psi\rangle=T(a)\int\psi(x)|x\rangle\,dx
=\int\psi(x)T(a)|x\rangle\,dx
\]
\[
\int\psi(x)T(a)|x\rangle\,dx
=\int\psi(x)|x+a\rangle\,dx
=\int\psi(x-a)|x\rangle\,dx
\]
又由 \(\psi_a(x)\) 的物理意义，\side{态与表象}{值得注意的是，\(\langle x|T(a)|\psi\rangle=[T(a)\psi](x)=\psi(x-a)\) 中，第一个 \(T(a)\) 指的是抽象量子算符，作用在量子态上；而第二个 \(T(a)\) 是抽象量子算符在坐标表象下的表示，作用在坐标表象波函数上. }
\[
\int\psi(x)T(a)|x\rangle\,dx=\int\psi_a(x)|x\rangle\,dx
\]
有
\[
\int\psi_a(x)|x\rangle\,dx=\int\psi(x-a)|x\rangle\,dx
\]
即
\[
\resultbox{\psi_a(x)=\psi(x-a)}
\]
现在，刚刚的结论严谨了.

\pondering{平移算子必须有哪些性质？}\label{pondering:properties}

先考察对一个态 \(|\psi\rangle\) 做两次平移操作：
\[
\begin{aligned}
T(a)T(b)\int\psi(x)|x\rangle\,dx
&=\int\psi(x)T(a)T(b)|x\rangle\,dx\\
&=\int\psi(x)T(a)|x+b\rangle\,dx\\
&=\int\psi(x)|x+a+b\rangle\,dx\\
&=\int\psi(x-a-b)|x\rangle\,dx
\end{aligned}
\]
即有
\[
\langle x|T(a)T(b)|\psi\rangle=\psi(x-a-b)
\]
可推出
\[
\resultbox{T(a)T(b)=T(a+b)=T(b)T(a)}
\]\side{复合律}{两次平移的总位移为 \(a+b\). 交换两次操作的顺序，总位移相同. }

由平移的物理意义知，平移 \(0\) 长度即不做变换，有
\[
T(0)=I
\]
那么有
\[
T(a)T(-a)=T(a-a)=T(0)=I
\]
\[
T(a)=T^{-1}(-a)
\]

接下来考虑这样一个情况，我们有一个待测态 \(\psi\) 和一个参考态 \(\phi\)，我们关心 \(\psi\) 在 \(\phi\) 上的投影 \(\langle\phi|\psi\rangle\). 现在我们同步搬动 \(\psi\) 和 \(\phi\)，它们的内积会改变吗？
\[
\begin{aligned}
\langle\phi|T^\dagger(a)T(a)|\psi\rangle
&=\langle\phi|T^\dagger(a)\int|x\rangle\langle x|\,dx\,T(a)|\psi\rangle\\ &=\int\phi^*(x-a)\psi(x-a)\,dx\\
&=\int\phi^*(x)\psi(x)\,dx=\langle\phi|\psi\rangle
\end{aligned}
\]

结论：内积不变. 有
\[
T^\dagger(a)T(a)=I
\]
由于 \(T(a)\) 已有逆算符，立即得到 \(T^\dagger(a)=T(a)^{-1}=T(-a)\)，所以 \(T(a)T^\dagger(a)=I\)，即
\[
\resultbox{T^\dagger(a)T(a)=I=T(a)T^\dagger(a)}
\]
故 \(T(a)\) 是幺正算符.

最后，我们讨论如果平移很小一段距离，平移算子的形式是怎样的？由物理意义显然：
\[
\lim_{a\to0}T(a)=I
\]
这对后续我们讨论「平移很小一段距离」很有用.

事实上，算子的极限的确切定义是一个很有价值的数学问题. 一种常见的定义（强连续）是：
\[
\forall|\psi\rangle,\quad\lim_{a\to0}\left\|(T(a)-I)|\psi\rangle\right\|=0
\]
\side{强连续性}{严格来说，由于 \(|x\rangle\) 是广义本征态，上式无法直接证明. 但若对一维无自旋粒子，其 Hilbert 空间在位置表象下可实现为 \(L^2(\mathbb{R})\)，则在此空间上，该连续性在数学上是可以被严格证明的. }

但注意，即使接受了 \(\lim_{a\to0}T(a)=I\)，也不自动保证
\[
\lim_{a\to0}\frac{T(a)-I}{a}|\psi\rangle
\]
对任意 \(|\psi\rangle\) 均存在.

\pondering{什么是生成元？为什么要研究无穷小变换？}

生成元和有限变换的关系，可以类比为微分和积分的关系. 现在我们有了平移操作 \(T(a)\)，自然考虑 \(T(a)\) 作用在量子态上时，在起步阶段（\(a=0\)），态平移的有多快. 直观起见，我们考察坐标表象波函数对 \(a\) 的变化率，由导数定义
\[
\left.\frac{\partial\psi_a(x)}{\partial a}\right|_{a=0}
=\lim_{\varepsilon\to0}\frac{\psi(x-\varepsilon)-\psi(x)}{\varepsilon}
=-\frac{d\psi(x)}{dx}
\]
我们发现这就是 \(\psi(x)\) 对 \(x\) 的导数！这是自然的，因为若我们以 \(\psi(x)\) 波包为参考系，就会看到坐标轴在移动，也即对 \(x\) 求导. 负号的差异源于主动视角和被动视角的转换.

若我们换用一个表象无关的视角呢？已知 \(\psi_a(x)=\langle x|T(a)|\psi\rangle\)，\(\psi(x)=\langle x|\psi\rangle\)，\(\psi(x-\varepsilon)=\langle x-\varepsilon|\psi\rangle=\langle x|T^\dagger(-\varepsilon)|\psi\rangle=\langle x|T(\varepsilon)|\psi\rangle\)，代入坐标表象下的导数关系：
\[
\left.\frac{\partial\langle x|T(a)|\psi\rangle}{\partial a}\right|_{a=0}
=\lim_{\varepsilon\to0}\frac{\langle x|T(\varepsilon)|\psi\rangle-\langle x|\psi\rangle}{\varepsilon}
\]
\(\langle x|\) 显然和 \(a\) 无关，且由基底的完备性：
\side{导数存在性}{正如之前提醒过的，我们取使这个极限存在的态 \(|\psi\rangle\).}
\[
\left.\frac{\partial T(a)|\psi\rangle}{\partial a}\right|_{a=0}
=\lim_{\varepsilon\to0}\frac{T(\varepsilon)|\psi\rangle-|\psi\rangle}{\varepsilon}
\]
现在我们在两侧乘上 \(i\) \side{生成元的自伴性}{为什么要乘 \(i\)？为了保证生成元自伴. 回忆 \(T^\dagger(a)=T(-a)\). 形式上，在 \(a=0\) 处对两边求导，右边因 \(-a\) 多出一个负号，于是 \(T'^\dagger(0)=-T'(0)\). 再看 \(iT'(0)\)：取伴随时，\(i\) 也变号，两个负号恰好抵消.}
\[
i\left.\frac{\partial T(a)|\psi\rangle}{\partial a}\right|_{a=0}
=i\lim_{\varepsilon\to0}\frac{T(\varepsilon)|\psi\rangle-|\psi\rangle}{\varepsilon}
\]
现在令 \(K=i\left.\frac{\partial T(a)}{\partial a}\right|_{a=0}\)，形式上地，\(K\) 可被理解为 \(T(a)\) 对 \(a\) 的导数在 \(a=0\) 处的取值.\side{生成元的导数定义}{事实上，通过严格的数学方法，可以证明：对使上面差商极限存在的态 \(|\psi\rangle\)，该极限就是 \(-iK|\psi\rangle\).} 则有
\[
K|\psi\rangle=i\lim_{\varepsilon\to0}\frac{T(\varepsilon)|\psi\rangle-|\psi\rangle}{\varepsilon}
\]
使用 Peano 余项将极限误差吸入小量并移项，有
\[
T(\varepsilon)|\psi\rangle=|\psi\rangle-i\varepsilon K|\psi\rangle+o(\varepsilon)
\]\side{Taylor 展开}{这个式子在形式上接近 Taylor 展开，但事实上我们不保证 \(T(a)|\psi\rangle\) 有足够高阶的导数.}

我们将 \(K\) 称为 \(T(a)\) 的生成元. 现在我们可以接受：生成元是算子的导数. 现在一个很棒的问题是，为什么我们只研究 \(a=0\) 处的导数而非所有 \(a\) 处的导数？回忆叠加性
\[
T(a)T(b)=T(a+b)
\]
若要研究 \(T(a+da)\)，仅需研究 \(T(da)T(a)\)，即在 \(|x+a\rangle\) 上研究 \(T(0+da)\) 的作用. 这也暗示了我们可以用无穷多个无穷小的平移叠加成一个有限平移.

\pondering{如何用生成元描述有限平移？为什么平移生成元就是动量？}

有了关于导数的信息，如何构造出原「函数」？答案显然是「积分」，或更一般地，求解微分方程.

由平移的叠加性和导数的定义，
\[
\begin{aligned}
\frac{d}{da}T(a)|\psi\rangle
&=\lim_{\varepsilon\to0}\frac{T(a+\varepsilon)|\psi\rangle-T(a)|\psi\rangle}{\varepsilon}\\
&=\lim_{\varepsilon\to0}\frac{T(\varepsilon)-I}{\varepsilon}T(a)|\psi\rangle\\
&=-iKT(a)|\psi\rangle
\end{aligned}
\]
于是事实上我们得到微分方程
\[
\frac{d}{da}T(a)|\psi\rangle=-iKT(a)|\psi\rangle
\]
结合 \(T(0)=I\) 的初始条件，我们立即得到
\[
T(a)=e^{-iaK}
\]
这就是用生成元描述连续变换的方法. 

值得注意的是，不仅仅对平移算子，对任意满足 \(T(a)T(b)=T(a+b)\) 的幺正连续算子，必有指数表征 \(T(a)=e^{-iaK}\)，\(K\) 为 \(T\) 的自伴生成元. 其原因我们已经简单证明过.\side{Stone 定理}{严格的数学讨论涉及 Stone 定理.}

接下来我们来讨论 \(K\) 算子对应的到底是哪个物理量. 注意到一个有趣的数学事实，考察 \(K\) 在坐表表象下的表示：
\[
\langle x|K|\psi\rangle=-i\frac{d}{dx}\psi(x)
\]
\[
K=-i\frac{d}{dx}
\]
求它的本征波函数
\[
-i\frac{d}{dx}\psi_k(x)=k\psi_k(x)
\]
得
\[
\psi_k(x)=Ce^{ikx}
\]
显然，本征波函数的形式是单色平面波，而 \(k\) 是波数.

现在我们知道了，平移生成元 \(K\) 是波数 \(k\) 的算子. 由 de Broglie 关系，我们可以定义「波函数的动量」
\[
p=\hbar k
\]
和对应的动量算子
\[
P=\hbar K
\]
然后我们就能得到一个早已熟悉的关系：
\[
P=-i\hbar\frac{\partial}{\partial x}
\]

不过，这里有一个很微妙的区别：我们定义了和位置表象下单色平面波函数对应的动量 \(p\)，但没有保证它就是波函数描述的粒子的动量（若它描述的确实是粒子的话）. 也就是讲，\(p=\hbar k\) 和 \(p=mv\) 定义的动量是否一致，仍有待商榷. \side{自由粒子}{对自由粒子的波函数，二者在数值上相等.}
事实上，它们在物理意义上是不同的，二者在什么条件下数值上相等，要考虑具体的动力学系统后讨论，这就是后话了.

\section{时间演化}

\pondering{如何描述时间演化？}

研究完了量子态在尺子上怎么移动，现在我们来研究它怎么随着钟表演化. ​	设想我们在 \(t_0\) 时刻制备了一个粒子，其态为 \(|\psi(t_0)\rangle\). 现在的问题是，我们手里只有 \(|\psi(t_0)\rangle\)，如何得到 \(|\psi(t)\rangle\)？

​我们灵光乍现，发现可以照抄处理空间平移的思路，给这个动作定义一个算子，记为 \(U(t)\)，使得
\[
|\psi(t_0+t)\rangle=U(t)|\psi(t_0)\rangle
\]
现在我们可以说「我们定义的是 \(t=0\) 时刻附近的时间演化，然后根据叠加性 \(U(t+dt)=U(dt)U(t)\)，仅需处理 \(U(t)\) 在 \(t=0\) 的导数即可…」

但是仔细想想，这样处理合理吗？我们先前对空间平移这样操作，是因为在不同空间位置做的空间平移操作是等效的. 而时间演化并不保证如此. 考虑给系统施加一个随时间变化的外场，事实上我们不能保证
\[
\left.U(dt)\right|_{t=t_1}
=\left.U(dt)\right|_{t=t_2}
\]
良好体现这一性质的方法也是显而易见的：把初始时间 \(t_0\) 也放进演化算子的自变量里去. \side{双参数演化算子族}{严格来讲，这里的 \(U(t,t_0)\) 是「双参数演化算子族」，和描述空间平移的「单参数幺正群」略有区别.} 同时为了记号简单，认为这里的 \(t\) 是演化末态的时间，而非时间的变化量，有
\[
|\psi(t)\rangle=U(t,t_0)|\psi(t_0)\rangle
\]

对封闭系统，我们接受量子力学的一条动力学公设\side{量子力学公设}{该公设得到了大量实验的证实.}：时间演化由线性幺正算子描述. 也就是说，若让两个态在相同的条件下演化，它们的内积保持不变.

​不过，无论如何，从 \(t_0\) 等到 \(t_1\)，再从 \(t_1\) 等到 \(t_2\)，总归是从 \(t_0\) 等到了 \(t_2\). 于是有
\[
U(t_2,t_1)U(t_1,t_0)=U(t_2,t_0)
\]
同时
\[
U(t_0,t_0)=I
\]
和逆演化
\[
U^\dagger(t,t_0)=U^{-1}(t,t_0)=U(t_0,t)
​
\]

之前讲过，满足复合律的幺正连续算子，都可以用取无穷小变换后构造 \(e\) 指数的方法研究. 不过在此之前，让我们回忆这个问题：在初见量子力学的时候，是用什么工具描述波函数随时间演化的？答案是：Schrödinger 方程. 既然时间演化算子可以描述量子态的时间演化，Schrödinger 方程也起到了同样的作用，那么二者应当是等价的. 接下来我们就研究这个问题.

\pondering{能从时间演化推出薛定谔方程吗？}

​先考察时刻 \(t\)，取一个态 \(|\psi(t)\rangle\). 仿照空间平移时定义 \(K\) 的做法，我们考察 \(U(t+\varepsilon,t)|\psi(t)\rangle\) 在 \(\varepsilon=0\) 处的变化率. 定义哈密顿算子 \(H(t)\)：\side{差商极限存在性}{和之前一样，这里我们默认 \(|\psi(t)\rangle\) 是使得差商极限存在的态.}
\[
H(t)|\psi(t)\rangle
=i\hbar\lim_{\varepsilon\to0}
\frac{U(t+\varepsilon,t)|\psi(t)\rangle-|\psi(t)\rangle}{\varepsilon}
\]
\side{凑能量量纲}{为什么这次多乘了一个 \(\hbar\)？量子态对时间求导，得到的量纲是时间的倒数，乘上 \(\hbar\)，恰好有能量的量纲. 至于它和真实的系统能量有什么物理联系？这是个好问题，马上就会解答.}
分子里的第一项恰好就是 \(|\psi(t+\varepsilon)\rangle\). 有
\[
\begin{aligned}
H(t)|\psi(t)\rangle
&=i\hbar\lim_{\varepsilon\to0}
\frac{|\psi(t+\varepsilon)\rangle-|\psi(t)\rangle}{\varepsilon}\\
&=i\hbar\frac{d}{dt}|\psi(t)\rangle
\end{aligned}
\]
我们就得到了 Schrödinger 方程的一般形式
\[
i\hbar\frac{d}{dt}|\psi(t)\rangle=H(t)|\psi(t)\rangle
\]

显而易见地，我们可以质疑：这只是从时间演化算子的幺正连续叠加性推出的结论，没有使用任何和系统相关的性质，那岂不是我们连粒子在弹簧上还是在无限深势阱里都不需要知道，就可以推演粒子的演化了？显然不可能. 如果还记得之前关于为何 \(U(t, t_0)\) 需要两个参量的讨论，就会发现关键在于，系统的信息体现在 \(H(t)\) 里. 

举个例子，若给定
\[
H(t)=\frac{P^2}{2m}+V(X,t)
\]
再将方程两侧投影到 \(\langle x|\) 上，用之前得到的 \(P=-i\hbar\,\partial_x\)，才有
\[
i\hbar\frac{\partial\psi(x,t)}{\partial t}
=-\frac{\hbar^2}{2m}\frac{\partial^2\psi(x,t)}{\partial x^2}
+V(x,t)\psi(x,t)
\]
这就是我们非常熟悉的形式了. 系统的关键信息由 \(H(t)=\frac{P^2}{2m}+V(X,t)\) 给出，时间演化的复合律本身并没有告诉我们这件事.

\pondering{哈密顿量有什么物理意义？}

​我们已经通过凑量纲的方式让 \(H(t)\) 拥有了能量量纲，接下来我们思考两个问题：这个东西有什么用？它代表的能量到底是个啥？

​先研究第一个问题，考虑 \(H\) 不显含时间的情况，即等同样长的时间，演化算子就相同，可以记 \(\tau=t-t_0\)，把 \(U(t,t_0)\) 简写为 \(U(\tau)\). 和空间平移时逐项比较，有
\[
T(a)=e^{-iaP/\hbar},\qquad U(\tau)=e^{-i\tau H/\hbar}
\]
所以，知道 \(H\)，就知道从给定的初态出发，等任意长一段时间以后会得到什么态. 这就是「哈密顿量生成时间演化」的意思.

​再研究第二个问题，取 \(H\) 的一个本征态，暂且记为 \(\lambda\)，有
\[
H|\lambda\rangle=\lambda|\lambda\rangle
\]
把 \(U(\tau)=e^{-i\tau H/\hbar}\) 做泰勒展开，逐项代入上式以后求和回去，有
\[
U(\tau)|\lambda\rangle=e^{-i\lambda\tau/\hbar}|\lambda\rangle
\]
回忆之前我们从 \(e^{ikx}\) 中读出了波数 \(k\). 现在把上面的时间因子写成 \(e^{-i\omega\tau}\)，读出 \(\lambda=\hbar\omega\). 用 de Broglie 关系 \(E=\hbar\omega\). 在选定能量零点后\side{整体相位}{事实上，改变能量零点意味着给量子态加上一个整体相位.}，\(H\) 的本征值就是系统能量的概念，可以记为 \(E\)，本征态也就可以记成 \(|E\rangle\)，就有
\[
H|E\rangle=E|E\rangle
\]
这便是我们早就熟悉的定态薛定谔方程. 顺便，我们还得到了不显含时哈密顿量系统能量本征态的时间演化规律：乘上一个形如 \(e^{-iE\tau/\hbar}\) 相位因子，这也是早已熟悉的. \side{显含时哈密顿量}{对显含时哈密顿量的处理方法，将在量子动力学部分详细介绍.}

不过有一个问题值得一问：之前我们不能确定 \(P\) 对应的物理量就是由 \(mv\) 定义的动量，这里我们怎么能保证 \(H\) 对应的物理量是系统的能量呢？一个解释是：经典力学里，时间平移对称性对应的 Noether 守恒荷就是能量，而在量子力学里相关对称性可以类比. 另一个解释是：大量实验事实都证明了这一点. 但是从纯粹的连续变换的数学结构证明这个量的物理意义是不可能的.

\pondering{为什么时间是参数，位置是算符？}

\side{时间算符}{如果我们直接把时间定义成算符，假设它和 \(H\) 拥有类似 \(X\) 和 \(P\) 的正则对易关系，会出现一些不好的下界发散问题. 详见 Pauli 时间算符论证，这里暂不展开.}
看位置表象波函数
\[
\psi(x,t)=\langle x|\psi(t)\rangle
\]
虽然 \(x,t\) 并排写在同一个括号里，右侧却告诉我们，它们进入这个式子的方式很不一样. 给定 \(t\)，我们取出此刻的态 \(|\psi(t)\rangle\). 给定 \(x\)，我们问这个态在位置本征态 \(|x\rangle\) 上的投影. 尺子上读到的 \(x\) 当然也是一个数，所谓「位置是算符」，说的是描述这项测量的可观测量 \(X\)，它满足 \(X|x\rangle=x|x\rangle\). 在这里讨论的非相对论量子力学里，实验室的钟给态标记时刻，直接读出就好了，不需要专门拿一个算符去测量. 所以我们把 \(t\) 当作参数.